% Three Independent Capabilities of Reflective Computation
% LICS submission — ACM SIGCONF format, double-blind
\documentclass[sigconf,screen,review,anonymous]{acmart}

% Lean code listings
\usepackage{listings}
\lstset{
  basicstyle=\ttfamily\small,
  columns=flexible,
  breaklines=true,
  frame=none,
  xleftmargin=1em,
}

% Theorems
\newtheorem{observation}{Observation}

% Shortcuts
\newcommand{\dotop}{\cdot}
\newcommand{\FRM}{\textsf{FRM}}
\newcommand{\DRM}{\textsf{DRM}}
\newcommand{\ICP}{\textsf{ICP}}
\newcommand{\Fin}{\mathrm{Fin}}
\newcommand{\core}{\mathrm{core}}
\newcommand{\Nat}{\mathbb{N}}

\begin{document}

\title{Three Independent Capabilities of Reflective Computation in Finite Algebras}

\begin{CCSXML}
<ccs2012>
<concept>
<concept_id>10003752.10003790.10003795</concept_id>
<concept_desc>Theory of computation~Algebraic language theory</concept_desc>
<concept_significance>500</concept_significance>
</concept>
<concept>
<concept_id>10003752.10003790.10002990</concept_id>
<concept_desc>Theory of computation~Logic and verification</concept_desc>
<concept_significance>500</concept_significance>
</concept>
</ccs2012>
\end{CCSXML}

\ccsdesc[500]{Theory of computation~Algebraic language theory}
\ccsdesc[500]{Theory of computation~Logic and verification}

\keywords{finite algebra, reflective computation, independence, Lean~4, machine-checked proofs}

\begin{abstract}
We identify three capabilities of reflective computation — self-simulation, self-description, and evaluator internalization — and prove they are fully independent in finite extensional magmas: no capability implies any other. Each capability has a single-concept, presentation-independent definition. Self-simulation is the existence of a section-retraction pair. Self-description is the classifier dichotomy: every non-absorber element is either entirely boolean-valued or entirely non-boolean-valued on the core. Evaluator internalization is the Internal Composition Property: some element's action on core factors non-trivially through two others, one core-preserving.

We prove independence by exhibiting Lean-verified finite counterexamples in all four non-trivial directions ($N\!=\!8$ and $N\!=\!10$), exhibit a tight coexistence witness at $N\!=\!10$, prove the three-category decomposition is an algebraic invariant, and show the magma setting is necessary: no associative magma admits the classifier dichotomy (ruling out semigroups, monoids, and groups). All results are machine-checked in Lean~4, zero \texttt{sorry}.
\end{abstract}

\maketitle

% ══════════════════════════════════════════════════════════════
\section{Introduction}
% ══════════════════════════════════════════════════════════════

When a programmer writes a meta-circular evaluator, how many things are they doing? The standard answer is one: building a self-interpreter. The evaluator encodes programs as data, classifies expressions by type, and reduces them — all in a single recursive function. These mechanisms feel inseparable because they always appear together.

But they need not. A system can encode its own elements without classifying them (a G\"odel numbering has no type system). It can classify without evaluating (a type checker need not run programs). It can evaluate without classifying (an untyped reducer applies rules without sorting expressions into categories). Each of these is familiar in isolation. What is not obvious is whether the combinations are forced — whether, say, the ability to self-evaluate \emph{requires} the ability to classify, or merely co-occurs with it in practice.

This paper answers the question for finite algebras: the three mechanisms are fully independent. We exhibit concrete finite structures where any one is present and the others absent. The separation is not a matter of expressiveness or convenience — it is a theorem, machine-checked in Lean.

Specifically, we define three capabilities of a finite extensional magma~\cite{burris81} (a set with a binary operation whose left-regular representation is faithful):

\begin{itemize}
  \item \textbf{Self-simulation (S):} the existence of a section-retraction pair on the core — elements can encode and decode each other.
  \item \textbf{Self-description (D):} the classifier dichotomy — every non-absorber element is either entirely boolean-valued or entirely non-boolean-valued on the core, creating a clean partition into zeros, classifiers, and non-classifiers.
  \item \textbf{Evaluator internalization (H):} the Internal Composition Property (ICP) — some element's action on core factors non-trivially as $a \dotop x = c \dotop (b \dotop x)$ with $b$ core-preserving, $a \neq b \neq c$ pairwise distinct, and $a$ taking $\geq 2$ distinct values.
\end{itemize}

We prove these three capabilities are fully independent: no capability implies any other. Four Lean-verified finite counterexamples establish all four non-trivial directions (the remaining two are definitional: both D and H are formalized on top of S). A Lean-verified witness shows all three capabilities coexist at $N\!=\!10$.

The three-category decomposition induced by D (zeros, classifiers, non-classifiers) is proved to be an algebraic invariant: it is preserved by all magma isomorphisms that respect the absorbers. The ICP is proved logically equivalent to the standard operational axioms (Compose + Inert), giving evaluator internalization a single-concept characterization on par with S and D.

All results are formalized in Lean~4~\cite{moura21} and verified by \texttt{lake build}, zero \texttt{sorry}. The proofs divide into universal algebraic theorems (proved by pure equational reasoning, no \texttt{decide}) and model-specific verifications (proved by \texttt{native\_decide} on concrete Cayley tables).

\paragraph{Contributions.}
\begin{enumerate}
  \item An independence theorem: S, D, and H are pairwise independent, with Lean-verified counterexamples at $N\!=\!8$ and $N\!=\!10$ (Section~\ref{sec:independence}).
  \item A coexistence witness at $N\!=\!10$ (Section~\ref{sec:bound}).
  \item Universal algebraic obstructions: no $\DRM$ has a right identity or is associative, ruling out semigroups, monoids, and groups (Section~\ref{sec:universal}).
  \item A decomposition invariance theorem: the three-category partition is preserved by isomorphisms (Section~\ref{sec:invariance}).
  \item A characterization of H: the ICP is logically equivalent to Compose+Inert (Section~\ref{sec:icp}).
\end{enumerate}

% ══════════════════════════════════════════════════════════════
\section{Definitions}
\label{sec:definitions}
% ══════════════════════════════════════════════════════════════

Throughout, $(S, \dotop)$ is a finite magma on $\Fin(n) = \{0, 1, \ldots, n-1\}$ with binary operation $\dotop : \Fin(n) \times \Fin(n) \to \Fin(n)$.

\begin{definition}[Faithful Retract Magma]
\label{def:frm}
A \emph{Faithful Retract Magma} $(\FRM)$ is a tuple $(S, \dotop, z_1, z_2, s, r)$ where:
\begin{enumerate}
  \item $z_1, z_2 \in S$ are distinct \emph{left-absorbers}: $z_i \dotop x = z_i$ for all $x$.
  \item No other element is a left-absorber.
  \item \emph{Extensionality}: if $a \dotop x = b \dotop x$ for all $x$, then $a = b$.
  \item $s, r \in \core(S) := S \setminus \{z_1, z_2\}$ form a \emph{retraction pair}: $r \dotop (s \dotop x) = x$ and $s \dotop (r \dotop x) = x$ for all $x \in \core(S)$, with $r \dotop z_1 = z_1$ (anchoring the retraction to the zero structure).
\end{enumerate}
\end{definition}

The retraction pair provides encoding infrastructure: element $a$ can be represented as $s^k(z_1)$ for its index $k$, and $r$ peels one layer. This is the finite algebraic analog of a G\"odel numbering.

\begin{definition}[Classifier Dichotomy / Dichotomic Retract Magma]
\label{def:drm}
A \emph{Dichotomic Retract Magma} $(\DRM)$ extends $\FRM$ with:
\begin{enumerate}
  \item A \emph{classifier} $c \in \core(S)$ with $c \dotop x \in \{z_1, z_2\}$ for all $x$.
  \item The \emph{classifier dichotomy}: for every $y \in \core(S)$, either
    \begin{itemize}
      \item $y \dotop x \in \{z_1, z_2\}$ for all $x \in \core(S)$ \quad ($y$ is a \emph{classifier}), or
      \item $y \dotop x \notin \{z_1, z_2\}$ for all $x \in \core(S)$ \quad ($y$ is a \emph{non-classifier}).
    \end{itemize}
  \item At least one non-classifier exists. (Without this, every $\FRM$ where all non-absorbers happen to be classifiers would vacuously satisfy D. The non-degeneracy condition ensures the dichotomy is a genuine partition into two inhabited classes, not a trivial collapse.)
\end{enumerate}
\end{definition}

The dichotomy creates a three-category decomposition $S = Z \sqcup C \sqcup N$ where $Z = \{z_1, z_2\}$ (absorbers), $C$ (classifiers), and $N$ (non-classifiers). No element can be partially classifying. This is a finite algebraic stratification: truth-producing elements are cleanly separated from computation-producing elements.

\begin{definition}[Internal Composition Property]
\label{def:icp}
A magma $(S, \dotop, z_1, z_2)$ satisfies the \emph{Internal Composition Property} $(\ICP)$ when:
\[
\exists\, a, b, c \in \core(S),\ \text{pairwise distinct, such that:}
\]
\begin{enumerate}
  \item \emph{Core-preservation}: $b \dotop x \notin \{z_1, z_2\}$ for all $x \in \core(S)$.
  \item \emph{Factorization}: $a \dotop x = c \dotop (b \dotop x)$ for all $x \in \core(S)$.
  \item \emph{Non-triviality}: $|\{a \dotop x : x \in \core(S)\}| \geq 2$.
\end{enumerate}
\end{definition}

ICP says: the left-regular representation admits a non-trivial factorization on core. One element's action decomposes as ``store, then process'' — a single fragment of internal composition. This is the evaluator internalization core: the algebra contains an element representing a composite of two of its own operations.

\begin{definition}[Capability S]
An $\FRM$ has capability $S$ (self-simulation) by definition.
\end{definition}

\begin{definition}[Capability D]
A $\DRM$ has capability $D$ (self-description) by definition.
\end{definition}

\begin{definition}[Capability H]
An $\FRM$ has capability $H$ (evaluator internalization) if it satisfies the $\ICP$.
\end{definition}

Note that D and H are both defined on top of $\FRM$, so both imply S. The four non-trivial independence directions are S $\not\Rightarrow$ D, D $\not\Rightarrow$ H, H $\not\Rightarrow$ D, and S $\not\Rightarrow$ H.

% ══════════════════════════════════════════════════════════════
\section{Universal Theorems}
\label{sec:universal}
% ══════════════════════════════════════════════════════════════

All theorems in this section hold for \emph{every} $\DRM$, proved by pure algebraic reasoning (no \texttt{decide}). They are formalized in \texttt{CatKripkeWallMinimal.lean}, \texttt{NoCommutativity.lean}, and \texttt{Functoriality.lean}.

\begin{theorem}[Three-Category Decomposition]
\label{thm:three-cat}
In any $\DRM$, every element is a zero, a classifier, or a non-classifier. The classes are pairwise disjoint.
\end{theorem}

\begin{theorem}[Classifier Wall]
\label{thm:wall}
$C \cap N = \emptyset$: no element is partially classifying.
\end{theorem}

\begin{theorem}[Retraction Pair Placement]
\label{thm:placement}
The section $s$ and retraction $r$ are both non-classifiers: $s, r \in N$.
\end{theorem}

\begin{theorem}[Asymmetry]
\label{thm:asymmetry}
No extensional magma with two distinct left-absorbers is commutative.
\end{theorem}

\begin{theorem}[Cardinality Bounds]
\label{thm:card}
$|S| \geq 4$ for any $\DRM$, and $|S| \geq 5$ when $s \neq r$. Both bounds are tight.
\end{theorem}

\begin{theorem}[No Right Identity]
\label{thm:no-rid}
No $\DRM$ has a right identity element.
\end{theorem}

\begin{theorem}[No Associativity]
\label{thm:no-assoc}
No $\DRM$ is associative. In particular, no $\DRM$ is a semigroup, monoid, or group.
\end{theorem}

\begin{proof}[Proof sketch]
Assume associativity. Then $\text{cls} \dotop (a \dotop b) = (\text{cls} \dotop a) \dotop b = \text{cls} \dotop a$ for all $a, b$, since $\text{cls} \dotop a \in \{z_1, z_2\}$ (classifier), and absorbers absorb on the left. So the classifier absorbs on the right.

In particular, $\text{cls} \dotop z_1 = \text{cls} \dotop (\text{cls} \dotop b_1) = \text{cls} \dotop \text{cls}$ and $\text{cls} \dotop z_2 = \text{cls} \dotop (\text{cls} \dotop b_2) = \text{cls} \dotop \text{cls}$, where $b_1, b_2$ are witnesses to $\text{cls}$ hitting each absorber (which must exist, else extensionality forces $\text{cls} = z_i$). Meanwhile, for any core $x$: $\text{cls} \dotop x = \text{cls} \dotop (r \dotop (s \dotop x)) = \text{cls} \dotop r$, using the retraction pair and the right-absorption property. So $\text{cls}$ is constant at $c = \text{cls} \dotop r \in \{z_1, z_2\}$. By extensionality, $\text{cls} = z_i$ for some $i$, contradicting $\text{cls} \neq z_1, z_2$.

The full proof is algebraic (no \texttt{decide}) and uses neither the dichotomy nor the non-classifier existence axiom — only the classifier, the retraction pair, and extensionality. Lean file: \texttt{CatKripkeWallMinimal.lean}.\qed
\end{proof}

\subsection{Self-Simulation Injectivity}

The following results show that capability S alone constrains the algebra — even without D or H. They hold for any $\FRM$ that self-simulates (proved in \texttt{SelfSimulation.lean}). Self-simulation means: there exists a term $t$ in the free term algebra over $S$ such that evaluating $t$ on the encodings of any $a, b$ returns the atom $a \dotop b$.

\begin{theorem}[Partial Application Injectivity]
\label{thm:injectivity}
If an $\FRM$ self-simulates, the partial application map $a \mapsto \text{eval}(\text{App}(t, \text{rep}(a)))$ is injective.
\end{theorem}

\begin{corollary}
The encoding $\text{rep} : S \to \text{Term}(S)$ is injective, and equal partial applications imply identical rows.
\end{corollary}

% ══════════════════════════════════════════════════════════════
\section{Decomposition Invariance}
\label{sec:invariance}
% ══════════════════════════════════════════════════════════════

\begin{definition}[$\DRM$ Isomorphism]
A \emph{$\DRM$ isomorphism} $\varphi : M_1 \to M_2$ is a bijection $\varphi : \Fin(n) \to \Fin(n)$ satisfying:
\begin{itemize}
  \item $\varphi(a \dotop_1 b) = \varphi(a) \dotop_2 \varphi(b)$ for all $a, b$
  \item $\varphi(z_1^{(1)}) = z_1^{(2)}$ and $\varphi(z_2^{(1)}) = z_2^{(2)}$
\end{itemize}
\end{definition}

\begin{theorem}[Functoriality]
\label{thm:functoriality}
Any $\DRM$ isomorphism preserves the three-category decomposition: $\varphi$ maps zeros to zeros, classifiers to classifiers, and non-classifiers to non-classifiers.
\end{theorem}

\begin{proof}[Proof sketch]
\emph{Zeros}: if $a$ is a left-absorber, then $\varphi(a) \dotop_2 \varphi(x) = \varphi(a \dotop_1 x) = \varphi(a)$ for all $x$ in the image of $\varphi$; by surjectivity, $\varphi(a)$ is a left-absorber.

\emph{Classifiers}: if $y \dotop_1 x \in \{z_1, z_2\}$ for all core $x$, then $\varphi(y) \dotop_2 \varphi(x) = \varphi(y \dotop_1 x) \in \{z_1^{(2)}, z_2^{(2)}\}$ for all core $\varphi(x)$; by surjectivity and injectivity (to exclude absorber preimages), $\varphi(y)$ is a classifier.

\emph{Non-classifiers}: symmetric argument using injectivity to transport $\notin \{z_1, z_2\}$.

The full proof is algebraic (no \texttt{decide}) and formalized in \texttt{Functoriality.lean}.\qed
\end{proof}

% ══════════════════════════════════════════════════════════════
\section{Independence}
\label{sec:independence}
% ══════════════════════════════════════════════════════════════

\begin{theorem}[Full Independence]
\label{thm:independence}
No capability implies any other. Specifically:
\begin{enumerate}
  \item $S \not\Rightarrow D$: There exists an $\FRM$ of size $8$ with a classifier element that violates the classifier dichotomy.
  \item $D \not\Rightarrow H$: There exists a $\DRM$ of size $10$ where no triple satisfies the $\ICP$.
  \item $H \not\Rightarrow D$: There exists an $\FRM$ of size $10$ satisfying $\ICP$ that violates the classifier dichotomy.
  \item $S \not\Rightarrow H$: The minimal $\DRM$ witness ($N\!=\!4$) trivially lacks $\ICP$ since it has only 2 non-absorber elements but $\ICP$ requires 3. Lean file: \texttt{ICP.lean}, theorem \texttt{s\_not\_implies\_icp}.
\end{enumerate}
\end{theorem}

\begin{proof}
Each direction is proved by exhibiting a concrete Cayley table and verifying the claimed properties by \texttt{native\_decide} (or \texttt{decide} for $N\!=\!4, 8$). The tables are:

\paragraph{$S \not\Rightarrow D$ ($N\!=\!8$, the Countermodel).}
An $8 \times 8$ Cayley table defining an $\FRM$ with $Q = 2$, $E = 3$ (retraction pair), classifier at index~4, and two mixed elements (indices~5 and~7) whose rows on core contain both boolean and non-boolean values. Lean file: \texttt{Countermodel.lean}.

\paragraph{$D \not\Rightarrow H$ ($N\!=\!10$).}
A $10 \times 10$ Cayley table defining a $\DRM$ satisfying the classifier dichotomy. Lean proves that no triple $(a, b, c)$ of pairwise distinct non-absorbers satisfies $\ICP$, by exhaustive search via \texttt{native\_decide} over the full existential (\texttt{HasICP}). Lean files: \texttt{Countermodels10.lean} (DRM verification), \texttt{ICP.lean} (theorem \texttt{dNotH\_no\_icp}).

\paragraph{$H \not\Rightarrow D$ ($N\!=\!10$, diagonal).}
A $10 \times 10$ Cayley table defining an $\FRM$ with a retraction pair, Branch, Compose ($\eta \dotop x = \rho \dotop (g \dotop x)$), Inert ($g$ preserves core), and Y fixed-point — all evaluation machinery — yet element~5 has both boolean and non-boolean outputs on core, violating the dichotomy. Lean proves both that $\ICP$ holds (witnessed by $a\!=\!8, b\!=\!6, c\!=\!7$) and that the dichotomy fails. Lean file: \texttt{Countermodels10.lean}.

All tables are generated by Z3 SAT solver and independently verified. The Cayley tables are frozen in the repository (\texttt{counterexamples.json}) for reproducibility.\qed
\end{proof}

\paragraph{Interpretation.}
The $H \not\Rightarrow D$ direction is the most surprising: an algebra can have full evaluation machinery (encoding, conditional dispatch, internal composition, recursion) while having no clean separation between classifiers and non-classifiers. The classifier dichotomy is not a consequence of computational power — it is an independent organizational property. An algebra can \emph{evaluate} without \emph{understanding itself as evaluating}.

% ══════════════════════════════════════════════════════════════
\section{Coexistence}
\label{sec:bound}
% ══════════════════════════════════════════════════════════════

\begin{theorem}[Coexistence Witness]
\label{thm:bound}
The three capabilities S, D, and H are simultaneously satisfiable at $N\!=\!10$.
\end{theorem}

\begin{proof}
A concrete $10 \times 10$ Cayley table is exhibited that is simultaneously a $\DRM$ (Lean-verified: \texttt{witness10\_drm}) and satisfies $\ICP$ (Lean-verified: \texttt{w10\_has\_icp}). The witness uses 10 pairwise distinct elements in all roles: 2 absorbers, 2 retraction pair, 1 classifier, 3~$\ICP$ witnesses, and 2 enrichment elements ($f$ for Branch dispatch, $Y$ for recursion). Lean file: \texttt{Witness10.lean}.\qed
\end{proof}

\paragraph{Tightness.}
The $N\!=\!10$ bound is tight when all distinguished elements are required to be pairwise distinct (10 elements need $N \geq 10$ by counting; the Lean witness achieves it). Whether S+D+H can coexist at smaller $N$ with role overlap (e.g., an $\ICP$ witness coinciding with a retraction pair element) is open.

% ══════════════════════════════════════════════════════════════
\section{ICP Characterization}
\label{sec:icp}
% ══════════════════════════════════════════════════════════════

The standard formulation of evaluator internalization uses two operational axioms:
\begin{itemize}
  \item \emph{Compose}: $\eta \dotop x = \rho \dotop (g \dotop x)$ for all $x \in \core(S)$
  \item \emph{Inert}: $g \dotop x \notin \{z_1, z_2\}$ for all $x \in \core(S)$
\end{itemize}
with $\eta, g, \rho$ pairwise distinct non-absorbers and $\eta$ non-trivial ($\geq 2$ distinct values on core).

\begin{theorem}[$\ICP$ Equivalence]
\label{thm:icp-equiv}
For any finite magma on a 2-pointed set $(S, \dotop, z_1, z_2)$ of any size $n$:
\[
\ICP \iff \text{Compose} + \text{Inert} \text{ (with non-trivial composite)}.
\]
\end{theorem}

\begin{proof}
This is a \emph{universal} theorem, quantified over all $n$ and all binary operations. The witness triple $(a, b, c)$ of $\ICP$ maps to $(g, \rho, \eta) = (b, c, a)$ for Compose+Inert, and vice versa. The nine conditions (pairwise distinct, non-absorber, core-preservation, factorization, non-triviality) match exactly under this relabeling. The proof is pure equational logic — no \texttt{decide}, no \texttt{native\_decide}, no case analysis on $n$. Lean file: \texttt{ICP.lean}, theorem \texttt{icp\_iff\_composeInert}.\qed
\end{proof}

This equivalence gives H a single-concept definition: the left-regular representation contains a non-trivially composed element. The operational axioms (Compose, Inert) and the $\ICP$ are the same property under variable renaming. Additionally, $\ICP$ is invariant under $\FRM$ isomorphisms (algebraic proof, no \texttt{decide}; \texttt{ICP.lean}, theorem \texttt{FRMIso.preserves\_icp}).

\paragraph{Enrichments.}
Two additional axioms are used in the full reflective architecture but are not part of the $\ICP$ core:
\begin{itemize}
  \item \emph{Branch}: $\rho \dotop x = f \dotop x$ if $c \dotop x = z_1$, else $\rho \dotop x = g \dotop x$. This connects D (the classifier's judgment) to H (the evaluator's dispatch).
  \item \emph{Y} (fixed-point): $Y \dotop \rho = \rho \dotop (Y \dotop \rho)$ with $Y \dotop \rho \in \core(S)$. This crosses the decidability boundary from bounded to unbounded evaluation.
\end{itemize}
Both are irredundant (SAT-verified) but not part of the core independence result. They connect the capabilities and add computational power.

% ══════════════════════════════════════════════════════════════
\section{Categorical Context}
\label{sec:categorical}
% ══════════════════════════════════════════════════════════════

Each capability has connections to standard categorical vocabulary, of varying strength:

\paragraph{S = section-retraction pair (exact).} The retraction pair $(s, r)$ with $r \circ s = \mathrm{id}$ on core is a standard categorical construction~\cite{maclane71}. Definition~\ref{def:frm} is a section-retraction pair in the endomorphism category of the magma. The connection is formal: the Lean formalization uses exactly this structure.

\paragraph{D $\approx$ decidable subobject classifier (analogy).} In a topos, the subobject classifier $\Omega$ partitions morphisms into $\Omega$-valued and non-$\Omega$-valued~\cite{johnstone02}. The classifier dichotomy (Definition~\ref{def:drm}) is a finite analog: every non-absorber is either entirely $\{z_1, z_2\}$-valued or entirely non-$\{z_1, z_2\}$-valued on core. The connection is an analogy, not a formal functor: our dichotomy requires totality (no partial classifiers), a finiteness-dependent condition with no direct topos-theoretic counterpart. What is proved: the three-category decomposition and its invariance under isomorphisms (Theorem~\ref{thm:functoriality}). What is analogical: the identification of $C$ with the subobject classifier image.

\paragraph{H: categorical motivation (open).} The $\ICP$ was motivated by the notion of internal composition in monoidal closed categories, where the internal hom $[B, C]$ represents $\mathrm{Hom}(B, C)$ as an object. The $\ICP$ asks for a weaker condition: the left-regular representation contains at least one non-trivially composed element $L_a = L_c \circ L_b$. Full closure ($\{L_a\}$ closed under composition) fails for all tested models. Unlike S, the connection to internal hom is not formal — the $\ICP$ is a concrete algebraic property that we prove equivalent to Compose+Inert (Theorem~\ref{thm:icp-equiv}) and invariant under isomorphisms, but whether it corresponds to a standard categorical universal property is open.

% ══════════════════════════════════════════════════════════════
\section{Related Work}
\label{sec:related}
% ══════════════════════════════════════════════════════════════

\paragraph{Reflective towers.}
Smith's 3-Lisp~\cite{smith84} introduced the reflective tower — an infinite chain of meta-interpreters. Subsequent work~\cite{friedman84,danvy88} explored reification, intensions, and live meta-modification in reflective towers. Amin and Rompf~\cite{amin18} showed how to compile through the tower via stage polymorphism. All of these systems combine self-simulation, self-description, and evaluator internalization without separating them. Our contribution is the separation and the proof that it is genuine (no capability implies any other).

\paragraph{Self-interpretation.}
Mogensen~\cite{mogensen92} constructs self-interpreters for the $\lambda$-calculus. Brown and Palsberg~\cite{brown16} show that self-interpretation is possible in strongly-normalizing typed languages, contradicting earlier impossibility results. These works study self-interpretation (our S) but do not address the independence of S from D or H.

\paragraph{Subobject classifiers.}
The classifier dichotomy (Definition~\ref{def:drm}) is the finite analog of a decidable subobject classifier in topos theory~\cite{johnstone02,maclane92}. The three-category decomposition parallels the partition of morphisms into $\Omega$-valued and non-$\Omega$-valued in a topos, but our dichotomy requires \emph{totality} (no partial classifiers), which is a finiteness-dependent condition.

\paragraph{Finite model theory.}
Independence results via finite counterexamples are standard in finite model theory~\cite{libkin04}. Our contribution is not the technique but the specific separation: three independently-motivated capabilities of reflective computation, shown independent in an algebraic setting, with machine-checked proofs.

% ══════════════════════════════════════════════════════════════
\section{Discussion}
\label{sec:discussion}
% ══════════════════════════════════════════════════════════════

\paragraph{The classifier dichotomy is epistemic.}
The $H \not\Rightarrow D$ result (Section~\ref{sec:independence}) shows that evaluation machinery does not force clean classification. A system can evaluate without organizing its elements into coherent roles. The classifier dichotomy is an independent \emph{organizational} property — it makes the algebra legible, not more powerful. This distinction between computational power and structural legibility is, to our knowledge, new.

\paragraph{Why magmas?}
Magmas are extremely general — no associativity, no identity. A natural concern is that the independence result is an artifact of weak structure: in semigroups or monoids, the capabilities might collapse. Two Lean-proved results address this. Theorem~\ref{thm:no-rid} proves that no $\DRM$ has a right identity, ruling out monoids. Theorem~\ref{thm:no-assoc} proves the strictly stronger result: no $\DRM$ is associative, ruling out semigroups, monoids, and groups entirely. The proof (algebraic, no \texttt{decide}) shows that associativity forces the classifier to have a constant row, collapsing it to an absorber. Together, these results show that the magma setting is not conveniently weak — stronger algebraic structures cannot support capability D at all.

\paragraph{Scope and limitations.}
Whether analogs of the three-capability separation exist in other settings — $\lambda$-calculus, typed systems, infinite algebras — remains the primary open question. The categorical context (Section~\ref{sec:categorical}) suggests the concepts translate, but the no-associativity result (Theorem~\ref{thm:no-assoc}) means any such translation must work in settings without associativity or identity. No formal cross-setting result exists.

The ICP characterization of H (Section~\ref{sec:icp}) is a universal theorem: $\ICP \iff \text{Compose} + \text{Inert}$ for any magma on a 2-pointed set, proved by pure logic in Lean. ICP is also invariant under $\FRM$ isomorphisms (algebraic proof, no \texttt{decide}). However, while S and D have direct categorical analogs (section-retraction~\cite{maclane71}, decidable subobject classifier~\cite{johnstone02}), ICP is a fragment of internal composition — weaker than a full internal hom. A categorical characterization of H as a universal property remains open.

\paragraph{Realization freedom.}
Each capability admits multiple axiom forms (e.g., 6 satisfiable Compose variants for H). The specific role inventory depends on the chosen presentation. Whether a canonical selection criterion exists among equivalent presentations is an open question explored in the companion artifact.

% ══════════════════════════════════════════════════════════════
\section{Conclusion}
% ══════════════════════════════════════════════════════════════

We have shown that three capabilities of reflective computation — self-simulation, self-description, and evaluator internalization — are independent in finite extensional magmas, each admitting a single-concept presentation-independent definition. The independence is machine-checked; the decomposition is invariant; the coexistence bound is tight; and the magma setting is necessary (no semigroup, monoid, or group can satisfy D). The classifier dichotomy is epistemic, not computational: an algebra can evaluate without understanding itself as evaluating.

\paragraph{Open problems.}
\begin{enumerate}
  \item \emph{Cross-formalism universality}: Does the three-capability separation appear in the $\lambda$-calculus, typed settings, or infinite algebras?
  \item \emph{Categorical characterization of H}: Is the ICP equivalent to a standard universal property (e.g., existence of a partial internal hom)?
  \item \emph{Tight bound with overlap}: $N\!=\!10$ is tight with all roles distinct. Can S+D+H coexist at smaller $N$ with role overlap (e.g., an ICP witness coinciding with a retraction pair element)?
\end{enumerate}

\paragraph{Artifact.}
All Lean files, counterexample tables, SAT reproduction scripts, and the 16-element reflective tower artifact are available at \texttt{[repository URL]}.

% ══════════════════════════════════════════════════════════════
% References
% ══════════════════════════════════════════════════════════════

\bibliographystyle{ACM-Reference-Format}
\begin{thebibliography}{10}

\bibitem{smith84}
Brian Cantwell Smith.
\newblock Reflection and semantics in {Lisp}.
\newblock In \emph{POPL}, 1984.

\bibitem{friedman84}
Daniel P. Friedman and Mitchell Wand.
\newblock Reification: Reflection without metaphysics.
\newblock In \emph{LFP}, 1984.

\bibitem{danvy88}
Olivier Danvy and Karoline Malmkj{\ae}r.
\newblock Intensions and extensions in a reflective tower.
\newblock In \emph{LFP}, 1988.

\bibitem{amin18}
Nada Amin and Tiark Rompf.
\newblock Collapsing towers of interpreters.
\newblock In \emph{POPL}, 2018.

\bibitem{mogensen92}
Torben Mogensen.
\newblock Efficient self-interpretation in lambda calculus.
\newblock \emph{Journal of Functional Programming}, 2(3), 1992.

\bibitem{brown16}
Matt Brown and Jens Palsberg.
\newblock Breaking through the normalization barrier: A self-interpreter for {F-omega}.
\newblock In \emph{POPL}, 2016.

\bibitem{maclane71}
Saunders Mac~Lane.
\newblock \emph{Categories for the Working Mathematician}.
\newblock Springer, 1971.

\bibitem{johnstone02}
Peter~T. Johnstone.
\newblock \emph{Sketches of an Elephant: A Topos Theory Compendium}, vol.~1.
\newblock Oxford University Press, 2002.

\bibitem{maclane92}
Saunders Mac~Lane and Ieke Moerdijk.
\newblock \emph{Sheaves in Geometry and Logic: A First Introduction to Topos Theory}.
\newblock Springer, 1992.

\bibitem{libkin04}
Leonid Libkin.
\newblock \emph{Elements of Finite Model Theory}.
\newblock Springer, 2004.

\bibitem{burris81}
Stanley Burris and H.~P. Sankappanavar.
\newblock \emph{A Course in Universal Algebra}.
\newblock Springer, 1981.

\bibitem{kripke75}
Saul Kripke.
\newblock Outline of a theory of truth.
\newblock \emph{Journal of Philosophy}, 72(19), 1975.

\bibitem{moura21}
Leonardo de~Moura and Sebastian Ullrich.
\newblock The {Lean}~4 theorem prover and programming language.
\newblock In \emph{CADE}, 2021.

\end{thebibliography}

\end{document}
